\section{硬件中断从电气请求变成 IDT 向量}

设备在状态变化后向中断控制器提出请求。控制器根据屏蔽、优先级和路由选择目标
CPU 与向量，处理器在指令边界接受中断，再通过 IDT 进入内核。软件读取设备
状态确认来源，处理完成后分别确认设备与控制器。

\begin{pdfdiagram}
  \node[os hardware] (device) {设备\\状态变化};
  \node[os state,right=of device] (controller) {中断控制器\\屏蔽、优先级、向量};
  \node[os block,right=of controller] (cpu) {CPU\\指令边界接受};
  \node[os state,below=of cpu] (idt) {IDT 门与入口栈};
  \node[os block,left=of idt] (handler) {驱动上半部\\确认并排队};
  \node[os hardware,left=of handler] (eoi) {设备完成与 EOI};
  \draw[os arrow] (device) -- (controller);
  \draw[os arrow] (controller) -- (cpu);
  \draw[os arrow] (cpu) -- (idt);
  \draw[os arrow] (idt) -- (handler);
  \draw[os arrow] (handler) -- (eoi);
\end{pdfdiagram}
\diagramnote{IRQ 不是普通函数调用；入口与返回跨越设备、控制器、CPU、栈和驱动五个状态机。}

若驱动先发 EOI 再清设备条件，电平触发中断可能立即重新进入；若只清设备而漏
EOI，控制器可能不再交付同优先级 IRQ。顺序来自控制器触发模式和设备规范。

\section{8259A 初始化命令字形成严格序列}

主从 PIC 上电后的向量与 x86 异常重叠。初始化先保存屏蔽，向 command 端口
写 ICW1，随后向 data 端口写向量基址、级联关系和 8086 模式，最后恢复只允许
已配置 IRQ 的屏蔽位。主片 IRQ2 连接从片，不能在使用从片时屏蔽。

向量基址要求低三位为零，因为控制器用它们编码 IRQ 号。选择
\code{0x20} 和 \code{0x28} 后，IRQ0--15 映射到 \code{0x20..0x2F}，
避开 CPU 的 \code{0x00..0x1F}。

\section{EOI 与虚假中断需要读取在服务状态}

普通主片 IRQ 向主片发送 EOI；从片 IRQ 先向从片发送，再向主片确认级联 IRQ2。
IRQ7 和 IRQ15 可能在处理器确认前撤销，形成虚假中断。驱动读取 ISR 判断该位
是否真正处于服务中，避免错误 EOI 破坏优先级状态。

PIC 是学习兼容 PC 中断的起点。APIC 把确认、优先级和路由扩展到多核，但仍
遵循“识别来源、处理条件、确认控制器”的基本模型。

\section{PIT 除数把硬件频率转换为系统 tick}

PIT 通道 0 输入频率约 1.193182 MHz。模式 2 周期发生器的输出频率为：
\[
  f_{\mathrm{tick}}=\frac{1\,193\,182}{\mathrm{divisor}}.
\]
若目标约 100 Hz，除数约 11932，必须分低字节和高字节写入。整数除法产生
微小误差，内核时间换算应保留实际编程频率，不把每 tick 永久假定为精确
10 ms。

中断处理器递增 64 位单调 tick，唤醒到期定时器并减少当前线程预算。墙钟时间、
时区和日历属于更高层，不能直接由 PIT tick 代替。

\section{中断延迟由关闭窗口和处理时间共同决定}

从设备请求到处理器入口的时间受 IF、优先级、当前指令、其他中断和临界区影响。
长时间关闭中断会增加键盘丢失、定时抖动和 I/O 完成延迟。系统需要测量最大
关闭窗口，而不只观察平均吞吐。

上半部只读取必须及时取走的数据、清除条件并排队工作。日志格式化、块复制和
文件系统操作移到下半部或工作线程。这样 IRQ 响应时间由短而可界定的路径决定。

\section{PS/2 控制器与键盘是两个串联状态机}

控制器状态端口报告输出缓冲、输入缓冲和错误；数据端口可能返回键盘扫描码，
也可能返回控制器命令响应。写命令前等待输入缓冲为空，读数据前等待输出缓冲
非空，所有等待都有超时。

扫描码集合定义按下、释放和扩展前缀。解析器维护前缀与修饰键状态，输出
\code{KeyEvent}；布局层再把物理键位映射为 Unicode 或控制动作。IRQ 只把
原始字节放入环形队列，避免在中断上下文执行复杂文本逻辑。

\section{环形缓冲需要明确满与空的表示}

固定容量环形队列保存 head、tail 和元素数组。若只比较 head 与 tail，必须
牺牲一个槽或增加 count/代际位来区分满与空。所有索引使用容量取模或条件回绕，
容量和索引类型要避免截断。

单生产者单消费者可以用原子索引与 acquire/release 顺序；IRQ 与单核线程也可
在极短窗口关闭本地中断。多生产者、跨 CPU 或可睡眠消费者需要更完整同步。
选择依据并发模型，不依据容器看起来是否简单。

\section{锁的正确性包含上下文与顺序}

自旋锁等待期间持续占用 CPU，只适合极短临界区；互斥锁可以把线程放入等待队列，
不能在中断上下文使用。获取会被本地 IRQ 再入的锁前，要保存并关闭本地中断，
释放后恢复原状态而不是无条件开启。

全局锁顺序防止线程 A 持有 X 等 Y、线程 B 持有 Y 等 X。文档为每把锁记录
保护对象、可用上下文、是否可睡眠和排序级别。调试构建记录持有者与递归获取，
在死锁变成整机停滞前报告。

\section{抢占状态与中断状态不是同一个开关}

关闭中断阻止本地硬件 IRQ，不阻止其他 CPU；禁用抢占阻止调度器切换当前线程，
不一定阻止 IRQ。持自旋锁时通常同时禁止抢占，IRQ-safe 锁还保存本地中断状态。

把这些机制混成一个布尔量会在多核阶段失效。内核使用带生命周期的 guard
对象表达恢复责任，即使早期实现底层仍是显式汇编。

\section{驱动把寄存器状态机封装成请求状态机}

块请求经历 queued、submitted、in-flight、completed 或 failed。写命令只从
submitted 推进到 in-flight，IRQ 或轮询确认才进入 completed。超时路径可能
需要复位控制器并使所有在途请求失败。

请求对象拥有缓冲区引用、方向、范围、截止时间和完成状态。调用者不能在完成前
释放缓冲区，驱动也不能在报告完成后继续 DMA。所有权与硬件完成语义必须同一
接口表达。

\section{从 Ring 0 进入 Ring 3 需要五项共同成立}

用户代码和数据页设置 U/S，GDT 提供 DPL=3 的代码与数据段，TSS.RSP0 指向
可信内核栈，用户栈已经映射，\code{iretq} 帧包含用户 SS、RSP、RFLAGS、
CS 和 RIP。任一项错误都会在边界处产生保护异常或页故障。

进入前清理不应泄漏的寄存器，限制 RFLAGS 中 IOPL 等敏感位。用户 RIP 与
RSP 必须是规范地址并位于对应用户 VMA，不能只依赖处理器在返回时发现错误。

\section{系统调用入口先离开不可信栈}

\code{syscall} 不会像中断门一样自动切换栈。入口通过受控 GS 基址取得当前
线程内核栈，保存用户 RSP，再切换到内核栈和统一帧。RCX 与 R11 按指令语义
保存返回 RIP 与 RFLAGS，属于固定破坏寄存器。

\code{sysret} 对非规范地址等边界敏感，初期可以使用更统一的 \code{iretq}
返回。无论选择哪种，系统调用 ABI 都固定编号、参数寄存器、返回值和错误表示，
用户库与内核共享版本化定义。

\section{用户内存复制是可能失败的内核操作}

用户给出的地址可能未映射、只读、跨越页面或在并发线程中变化。内核先检查
\([address,address+length)\) 加法，再逐页验证用户权限，并在受控故障恢复
区域复制。验证与使用必须受地址空间锁或页面固定保护。

字符串接口还要限制最大长度并保证终止。内核不能先用普通
\code{strlen} 扫描未知用户地址，再检查范围；扫描本身已经可能越界故障。

\section{用户错误与内核错误采用不同终止边界}

用户非法指令、越权页访问和坏系统调用只终止当前进程并释放资源；内核在持锁
或修改全局结构时发生同类异常通常需要 panic。异常分发器根据 CPL、故障地址
和当前恢复上下文区分。

隔离验收不仅运行正常程序，还让用户代码执行 \code{out}、访问内核页、传入
跨页指针和制造除零。内核应持续运行其他进程，并留下稳定退出原因。
